\documentclass[11pt,a4paper]{article}

% =========================================================
% PAGE LAYOUT
% =========================================================

\usepackage[
    a4paper,
    left=2cm,
    right=2cm,
    top=1.8cm,
    bottom=1.8cm
]{geometry}

% =========================================================
% ESSENTIAL PACKAGES
% =========================================================

\usepackage[T1]{fontenc}
\usepackage[utf8]{inputenc}
\usepackage{lmodern}

\usepackage{array}
\usepackage{booktabs}
\usepackage{tabularx}
\usepackage{longtable}
\usepackage{enumitem}

\usepackage{graphicx}
\usepackage{caption}
\usepackage{xcolor}
\usepackage{hyperref}
\usepackage{placeins}
\usepackage{tikz}
\usetikzlibrary{positioning,arrows.meta}

\usepackage{listings}
\usepackage{fancyvrb}

% =========================================================
% GENERAL FORMATTING
% =========================================================

\setlength{\parindent}{0pt}
\setlength{\parskip}{4pt}

\raggedbottom

\renewcommand{\arraystretch}{1.15}
\setlength{\tabcolsep}{5pt}

\graphicspath{{images/}{./}}

\newcolumntype{L}{>{\raggedright\arraybackslash}X}

\newcolumntype{P}[1]{%
    >{\raggedright\arraybackslash}p{#1}
}

% =========================================================
% LIST FORMATTING
% =========================================================

\setlist[itemize]{
    leftmargin=1.5em,
    itemsep=2pt,
    topsep=3pt
}

\setlist[enumerate]{
    leftmargin=1.5em,
    itemsep=2pt,
    topsep=3pt
}

% =========================================================
% HYPERLINKS
% =========================================================

\hypersetup{
    colorlinks=true,
    linkcolor=blue,
    urlcolor=blue,
    citecolor=blue,
    bookmarksnumbered=true,
    pdftitle={Sensor Metadata Management Platform - Phase 3},
    pdfauthor={Chrystaleni Argyrou, Antri Kelli, Argyris Fellas}
}

% =========================================================
% FIGURE FORMATTING
% =========================================================

\captionsetup{
    font=small,
    labelfont=bf,
    justification=centering,
    skip=5pt
}

\setlength{\textfloatsep}{10pt}
\setlength{\floatsep}{8pt}
\setlength{\intextsep}{8pt}

\renewcommand{\topfraction}{0.9}
\renewcommand{\bottomfraction}{0.8}
\renewcommand{\textfraction}{0.08}
\renewcommand{\floatpagefraction}{0.7}

% =========================================================
% CODE LISTINGS
% =========================================================

\definecolor{codebackground}{RGB}{247,247,247}
\definecolor{codeborder}{RGB}{220,220,220}
\definecolor{codecomment}{RGB}{70,110,70}

\lstdefinestyle{platformcode}{
    basicstyle=\ttfamily\footnotesize,
    backgroundcolor=\color{codebackground},
    frame=single,
    rulecolor=\color{codeborder},
    breaklines=true,
    breakatwhitespace=false,
    columns=fullflexible,
    keepspaces=true,
    showstringspaces=false,
    numbers=none,
    tabsize=2,
    xleftmargin=5pt,
    xrightmargin=5pt,
    framexleftmargin=5pt,
    framexrightmargin=5pt,
    aboveskip=8pt,
    belowskip=8pt
}

\lstset{style=platformcode}

% =========================================================
% DOCUMENT
% =========================================================

\begin{document}

% =========================================================
% TITLE PAGE
% =========================================================

\begin{titlepage}
    \centering

    \vspace*{2cm}

    {\Large University of Cyprus\par}

    \vspace{0.5cm}

    {\large Department of Electrical and Computer Engineering\par}

    \vspace{2.5cm}

    {\Huge\bfseries Sensor Metadata Management Platform\par}

    \vspace{0.8cm}

    {\Large Phase 3 --- Software Implementation and Integration\par}

    \vspace{0.6cm}

    {\large In collaboration with the KIOS Research and Innovation Center of Excellence\par}

    \vfill

    {\large HMY317 --- Computing Technology\par}

    \vspace{1cm}

    {\large
    Team Members:\\[0.3cm]
    Chrystaleni Argyrou UC1366251\\
    Antri Kelli UC1071092\\
    Argyris Fellas UC1369020
    \par}

    \vspace{1.8cm}

    {\large November 2026\par}

\end{titlepage}

% =========================================================
% TABLE OF CONTENTS
% =========================================================

\tableofcontents

\newpage

% =========================================================
% 1. INTRODUCTION
% =========================================================

\section{Introduction}

\subsection{Purpose of Phase 3}

Phase 3 focuses on the construction, integration, testing, and
validation of the Sensor Metadata Management Platform developed
for the KIOS Research and Innovation Center of Excellence
(KIOS CoE).

The implementation builds upon the requirements identified
during Phase 1 and the architectural and object-oriented
design established during Phase 2.

The objective is to deliver an integrated web application
capable of managing metadata associated with sensors and
water distribution networks.

The platform provides centralized management of physical
devices, network elements, measurements, and controlled
reference lists.

It also supports Department-based authorization, external
system integration through a REST API, geographical
visualization, and Excel import/export operations.

\subsection{Implementation Objectives}

The principal implementation objectives are:

\begin{itemize}
    \item Implement persistent storage for the domain entities
    identified during requirements analysis.

    \item Develop REST API endpoints for metadata management
    and external system integration.

    \item Implement authentication and Department-based
    authorization.

    \item Develop a responsive web interface for managing
    and visualizing metadata.

    \item Support geographical information using PostGIS
    and GeoJSON.

    \item Implement Excel import functionality for both
    platform-specific and client-provided spreadsheets.

    \item Integrate the frontend, backend, and database
    into a deployable application.

    \item Validate the implementation through automated
    tests and integration checks.
\end{itemize}

% =========================================================
% 2. TECHNOLOGY SELECTION
% =========================================================

\section{Implementation Technologies}

\subsection{Technology Stack}

The implementation follows a three-tier architecture
consisting of a frontend application, backend services,
and a relational database.

Table~\ref{tab:phase3_technologies} summarizes the
technologies used during implementation.

\begin{table}[htbp]
\centering
\caption{Implementation Technology Stack}
\label{tab:phase3_technologies}
\small
\begin{tabularx}{\textwidth}{P{3.2cm} L}
\toprule
\textbf{Component} & \textbf{Technology} \\
\midrule
Backend &
Django 5.2 and Django REST Framework \\

Database &
PostgreSQL and PostGIS \\

Geospatial Support &
GeoDjango, GeoJSON, WGS84 (SRID 4326) \\

Frontend &
React 19 and Vite \\

Mapping &
Leaflet, MapLibre GL, and OpenFreeMap (OpenStreetMap data) \\

Backend Testing &
pytest \\

Frontend Testing &
Vitest, ESLint, and Playwright \\

Continuous Integration &
GitHub Actions \\

Deployment &
Docker Compose and nginx \\

API Documentation &
OpenAPI 3 and Swagger UI \\

Authentication &
Django sessions, CSRF protection, and API keys \\
\bottomrule
\end{tabularx}
\end{table}

\subsection{Backend Framework Decision}

During Phase 2, FastAPI was selected as the backend
framework.

However, following a request from KIOS CoE, the
backend implementation was changed to Django.

Django 5.2 was adopted as the main backend framework,
while Django REST Framework was selected for the
implementation of the REST API.

This change required adapting the original architectural
design to Django models, database migrations, serializers,
authentication mechanisms, and permission policies.

The change from FastAPI to Django represents a
client-requested architectural decision made after
the completion of Phase 2.

\subsection{Database and Geospatial Technologies}

PostgreSQL provides relational data storage, while
PostGIS extends the database with spatial capabilities.

GeoDjango is used to manage geographical fields
and spatial information.

The platform uses the WGS84 coordinate reference system,
identified by SRID 4326.

Supported geometry types include:

\begin{itemize}
    \item \texttt{Point}
    \item \texttt{LineString}
    \item \texttt{Polygon}
\end{itemize}

Geographical coordinates follow the GeoJSON convention:

\begin{lstlisting}
[longitude, latitude]
\end{lstlisting}

This representation allows geographical information
to be exchanged through the REST API and displayed
on the interactive map.

% =========================================================
% 3. SYSTEM ARCHITECTURE
% =========================================================

\section{System Architecture Implementation}

\subsection{Architectural Overview}

The implemented system consists of three principal
components:

\begin{enumerate}
    \item React frontend application.
    \item Django REST API backend.
    \item PostgreSQL/PostGIS database.
\end{enumerate}

The frontend communicates with the backend through
HTTP requests to the REST API.

The backend handles authentication, authorization,
validation, business logic, and database operations.

The database stores metadata records, organizational
relationships, and geographical information.

\subsection{Component Integration}

The principal communication flow is:

\begin{enumerate}
    \item A user interacts with the React interface.
    \item The frontend sends a request to the Django API.
    \item Django authenticates the request.
    \item The authorization policy evaluates the
    user's role and Department access.
    \item DRF serializers validate the request data.
    \item Django performs the corresponding database operation.
    \item The API returns a response to the frontend.
    \item React updates the displayed information.
\end{enumerate}

This architecture separates presentation logic
from backend business rules and persistent storage.

\subsection{Docker Deployment Architecture}

Docker Compose is used to coordinate three services:

\begin{itemize}
    \item \textbf{db:} PostgreSQL with PostGIS.
    \item \textbf{backend:} Django application and REST API.
    \item \textbf{frontend:} React application served through nginx.
\end{itemize}

The nginx configuration serves the frontend application
and forwards only the following paths to Django:

\begin{lstlisting}
/api/
/health/
/ready/
\end{lstlisting}

The health endpoints serve different purposes.

\begin{itemize}
    \item \texttt{/health/} provides a liveness check.
    \item \texttt{/ready/} verifies database connectivity
    and checks for pending migrations.
\end{itemize}

The Docker Compose configuration provides a local
deployment environment. Additional security configuration
would be required before exposing the application
to the public internet.

% =========================================================
% 4. DEVIATIONS FROM PHASE 2
% =========================================================

\section{Deviations from the Phase 2 Design}

The implementation follows the Phase~2 design. The three-tier modular
architecture, the Docker Compose deployment, the four roles with
multi-Department access, the HTTP~404/403 access rules, logical deletion,
the audit fields, the atomic Excel import, the read-only Department-scoped
API keys, and the separate Data Collection Configuration entity were
implemented as designed.

The deviations and extensions listed in
Table~\ref{tab:phase3_deviations} were introduced during implementation,
together with their justification.

\begin{longtable}{P{2.7cm} P{4.2cm} P{4.6cm} P{3.7cm}}
\caption{Deviations from the Phase 2 Design}
\label{tab:phase3_deviations}\\

\toprule
\textbf{Area} & \textbf{Phase 2 design} & \textbf{Implementation} & \textbf{Reason} \\
\midrule
\endfirsthead

\toprule
\textbf{Area} & \textbf{Phase 2 design} & \textbf{Implementation} & \textbf{Reason} \\
\midrule
\endhead

\bottomrule
\endfoot

Backend framework &
FastAPI with API routers, services, and a repository-based
Persistence Layer &
Django~5.2 and Django REST Framework. Models, serializers, and viewsets
replace the routers and repositories, and the Django ORM provides
persistence &
Request by KIOS CoE \\

API documentation &
Automatic FastAPI documentation at \texttt{/docs} &
OpenAPI~3 schema at \texttt{/api/v1/openapi.json} and Swagger UI at
\texttt{/api/v1/docs/}, both requiring authentication &
Consequence of the framework change \\

HTTP status codes &
400/422 for invalid input and 401 for authentication failures &
400 for invalid input. Unauthenticated session requests receive 403,
while invalid, expired, or revoked API keys receive 401 &
Django REST Framework conventions \\

Controlled lists &
One Controlled List Entry entity with categories, to which System
Administrators and authorized Department Administrators may add values &
Four separate lists (physical device types, network element types,
measurement types, and measurement units), owned by the Platform
Department and managed only by System Administrators &
Typed relations and stricter governance of the shared lists \\

API key expiration &
Expiration is optional &
Keys expire by default after 90 days, and the expiration date can
be set at creation &
Reduces the risk of long-lived credentials \\

Excel import format &
Excel import of metadata using the platform's workbook template &
The template is retained, and the client's workbook
(\texttt{EXCEL-317.xlsx}) is also recognized automatically &
The client's real data uses a different layout \\

Controlled-list values during import &
Values are added through list management &
Missing types and units found in a client workbook are created
automatically, but only when a System Administrator performs the import &
Avoids manual preparation of the lists before each import \\

Additional metadata fields &
Fields defined in the domain model &
Added \texttt{network\_area}, \texttt{elevation}, and \texttt{address}
(Network Element) and \texttt{communication\_parameters}
(Physical Device) &
These values exist in the client's workbook \\

\end{longtable}

% =========================================================
% 5. DATABASE IMPLEMENTATION
% =========================================================

\section{Database and Domain Model Implementation}

\subsection{Implemented Domain Entities}

The database implementation follows the domain model
defined during the previous project phases.

Table~\ref{tab:phase3_entities} summarizes the
principal entities.

\begin{longtable}{P{4.3cm} P{10.8cm}}
\caption{Implemented Domain Entities}
\label{tab:phase3_entities}\\

\toprule
\textbf{Entity} & \textbf{Description} \\
\midrule
\endfirsthead

\toprule
\textbf{Entity} & \textbf{Description} \\
\midrule
\endhead

\bottomrule
\endfoot

Organisation &
Represents an organization containing one or more
Departments. \\

Department &
Defines the organizational scope used for metadata
ownership and access control. \\

PhysicalDevice &
Represents a physical device associated with
measurement collection. \\

NetworkElement &
Represents an element of the monitored network,
including geographical and descriptive information. \\

Measurement &
Represents a measurement associated with a
PhysicalDevice and a NetworkElement. \\

ExternalIdentifier &
Stores external source/value identifiers used
to associate platform records with external systems. \\

DataCollectionConfiguration &
Stores configuration information associated
with a Measurement. \\

APIKey &
Represents a credential used by external systems
for read-only API access. \\

PhysicalDeviceType &
Controlled list of physical device types. \\

NetworkElementType &
Controlled list of network element types. \\

MeasurementType &
Controlled list of measurement types. \\

MeasurementUnit &
Controlled list of measurement units. \\

\end{longtable}

\subsection{Department Ownership}

Each business record belongs to a Department.

Each Department belongs to an Organisation.

Department ownership is used to restrict access
to records according to the authenticated user's
permissions.

Controlled reference lists are associated with
the Platform Department and are shared across
the platform.

Active entries in these lists must have
case-insensitively unique names.

\subsection{Soft Deletion}

The platform implements soft deletion for metadata
records.

Instead of permanently removing a record,
the system marks it as deleted.

Deleted records remain stored in the database
but are excluded from normal application queries.

This approach preserves stored information while
preventing deleted records from appearing in
ordinary user operations.

\subsection{Measurement Implementation}

The Measurement entity is associated with:

\begin{itemize}
    \item A PhysicalDevice.
    \item A NetworkElement.
    \item A MeasurementType.
    \item A MeasurementUnit.
\end{itemize}

Additional attributes include:

\begin{itemize}
    \item Description.
    \item Customer account.
    \item Meter account.
\end{itemize}

Measurement records may also be associated with
a separate DataCollectionConfiguration entity.

\subsection{Data Collection Configuration}

DataCollectionConfiguration is implemented as
a separate entity with a one-to-one relationship
to Measurement.

Database migration \texttt{0006} introduced this
structure and transferred existing configuration
data into the new representation.

The configuration can be accessed through:

\begin{lstlisting}
GET /api/v1/measurements/{id}/configuration/
\end{lstlisting}

If no configuration exists, the endpoint returns
HTTP 404.

The Measurement representation also includes
the configuration through the
\texttt{data\_collection\_config} field.

\subsection{Additional Metadata Fields}

Database migration \texttt{0007} introduced
additional metadata fields.

\textbf{NetworkElement}

\begin{itemize}
    \item \texttt{network\_area}
    \item \texttt{elevation}
    \item \texttt{address}
\end{itemize}

\textbf{PhysicalDevice}

\begin{itemize}
    \item \texttt{communication\_parameters}
    \item \texttt{location}
\end{itemize}

Communication parameters are represented using
JSON data.

The device location is represented as a
geographical Point.

\subsection{External Identifiers}

External identifiers consist of a source/value pair.

These identifiers support integration with external
systems and allow platform records to be associated
with identifiers originating outside the application.

The source/value combination is unique within
each Department.

% =========================================================
% 6. AUTHENTICATION AND AUTHORIZATION
% =========================================================

\section{Authentication and Authorization}

\subsection{Role-Based Access Control}

The platform implements four authorization roles:

\begin{enumerate}
    \item Viewer.
    \item Editor.
    \item Department Administrator.
    \item System Administrator.
\end{enumerate}

Access is granted only when the user has an
explicit platform role and the required
Department membership.

Django staff and superuser status do not
automatically grant platform permissions.

\subsection{Permission Matrix}

Table~\ref{tab:phase3_permissions} summarizes
the implemented access rules.

\begin{table}[htbp]
\centering
\caption{Role-Based Permission Matrix}
\label{tab:phase3_permissions}
\scriptsize
\begin{tabularx}{\textwidth}{
P{3.5cm}
P{2.2cm}
P{2.2cm}
P{2.5cm}
L
}
\toprule
\textbf{Operation} &
\textbf{Viewer} &
\textbf{Editor} &
\textbf{Dept. Admin} &
\textbf{System Admin} \\
\midrule

Read metadata &
Own Depts. &
Own Depts. &
Own Depts. &
All \\

Create/update metadata &
No &
Own Depts. &
Own Depts. &
All \\

Excel import &
No &
Own Depts. &
Own Depts. &
All \\

Read controlled lists &
Yes &
Yes &
Yes &
Yes \\

Manage controlled lists &
No &
No &
No &
Yes \\

Manage API keys &
No &
No &
Own Depts. &
All \\

Manage users &
No &
No &
Viewer/Editor &
All \\

Manage Departments &
No &
No &
No &
Yes \\

Manage Organisations &
No &
No &
No &
Yes \\

\bottomrule
\end{tabularx}
\end{table}

\subsection{Department Access Policy}

Department-based authorization is enforced
by the backend.

The default \texttt{NoAccessPolicy} denies
metadata access until the
\texttt{DepartmentAccessPolicy} is configured.

The implementation distinguishes between
two authorization failures:

\begin{itemize}
    \item \textbf{HTTP 404:} The requested record
    is outside the user's accessible scope.

    \item \textbf{HTTP 403:} The user can access
    the relevant scope but lacks permission
    to perform the requested operation.
\end{itemize}

This distinction prevents users from obtaining
information about records outside their
authorized Departments.

\subsection{Session Authentication}

The web application uses Django session
authentication.

Authentication is maintained through
session cookies and CSRF protection.

The implementation includes:

\begin{itemize}
    \item CSRF token validation.
    \item Throttling of unsuccessful login attempts.
    \item Session expiration checks.
    \item Password modification.
    \item Logout functionality.
\end{itemize}

Passwords and authentication tokens are not
stored in browser localStorage.

The password change endpoint is:

\begin{lstlisting}
POST /api/v1/auth/password/
\end{lstlisting}

After a successful password change,
the user remains authenticated.

\subsection{Authentication Context}

The frontend retrieves the current user's
authorization context through:

\begin{lstlisting}
GET /api/v1/auth/context/
\end{lstlisting}

The response provides information about:

\begin{itemize}
    \item The current platform role.
    \item Accessible Departments.
    \item Write permissions for each Department.
    \item Capability flags used by the interface.
\end{itemize}

The frontend uses this information to display
the appropriate workspace and settings.

However, all authorization decisions are
enforced by the backend.

\subsection{API Key Authentication}

External systems can authenticate using
API keys.

The request header follows the format:

\begin{lstlisting}
X-API-Key: smk_<prefix>.<secret>
\end{lstlisting}

API keys provide read-only access.

Requests attempting to modify data using
an API key return HTTP 403.

The implementation includes the following
security properties:

\begin{itemize}
    \item Random 256-bit secrets.
    \item SHA-256 digest storage.
    \item Constant-time digest comparison.
    \item One-time secret display during creation.
    \item Expiration, with a default validity
    period of 90 days.
    \item Revocation through DELETE requests.
    \item Preservation of the audit record
    after revocation.
\end{itemize}

API keys cannot be used to manage
other API keys or users.

% =========================================================
% 7. REST API IMPLEMENTATION
% =========================================================

\section{REST API Implementation}

\subsection{API Structure}

The REST API is exposed under:

\begin{lstlisting}
/api/v1/
\end{lstlisting}

The API supports metadata management,
authentication, organizational administration,
and Excel import operations.

\subsection{Metadata Endpoints}

Table~\ref{tab:phase3_api} summarizes the
principal API endpoints.

\begin{longtable}{P{6.3cm} P{8.8cm}}
\caption{REST API Endpoints}
\label{tab:phase3_api}\\

\toprule
\textbf{Endpoint} & \textbf{Functionality} \\
\midrule
\endfirsthead

\toprule
\textbf{Endpoint} & \textbf{Functionality} \\
\midrule
\endhead

\bottomrule
\endfoot

\texttt{physical-devices/} &
CRUD operations for physical devices. \\

\texttt{devices/} &
Alias for physical device operations. \\

\texttt{network-elements/} &
CRUD operations for network elements. \\

\texttt{measurements/} &
CRUD operations for measurements. \\

\texttt{physical-device-types/} &
Controlled list of physical device types. \\

\texttt{network-element-types/} &
Controlled list of network element types. \\

\texttt{measurement-types/} &
Controlled list of measurement types. \\

\texttt{measurement-units/} &
Controlled list of measurement units. \\

\texttt{measurements/\{id\}/configuration/} &
Retrieve measurement collection configuration. \\

\texttt{api-keys/} &
Create, retrieve, and revoke API keys. \\

\texttt{users/} &
User retrieval, creation, modification,
and deactivation. \\

\texttt{departments/} &
Retrieve Departments and create new Departments
with System Administrator permissions. \\

\texttt{organisations/} &
Retrieve Organisations and create new Organisations
with System Administrator permissions. \\

\texttt{imports/excel/} &
Import metadata from an Excel workbook. \\

\texttt{auth/context/} &
Retrieve authentication and authorization context. \\

\texttt{auth/password/} &
Change the authenticated user's password. \\

\end{longtable}

\subsection{Filtering and Pagination}

Metadata endpoints support:

\begin{itemize}
    \item Pagination.
    \item Search.
    \item Filtering.
    \item CRUD operations according to permissions.
\end{itemize}

Department access restrictions are applied
to API operations.

\subsection{User Management}

The user management API supports the following
operations:

\begin{lstlisting}
GET    /api/v1/users/
POST   /api/v1/users/
PATCH  /api/v1/users/{id}/
DELETE /api/v1/users/{id}/
\end{lstlisting}

DELETE deactivates the user rather than
permanently removing the account.

Department Administrators may manage
Viewer and Editor accounts within their
Departments.

System Administrators may manage all users.

\subsection{API Documentation}

The platform provides an OpenAPI 3 specification
through:

\begin{lstlisting}
GET /api/v1/openapi.json
\end{lstlisting}

Interactive Swagger documentation is
available at:

\begin{lstlisting}
GET /api/v1/docs/
\end{lstlisting}

Access to the documentation interface
requires authentication.

% =========================================================
% 8. EXCEL IMPORT AND EXPORT
% =========================================================

\section{Excel Import and Export Implementation}

\subsection{Import Functionality}

The platform supports Excel-based metadata
import through:

\begin{lstlisting}
POST /api/v1/imports/excel/
\end{lstlisting}

The endpoint accepts a multipart request
containing:

\begin{itemize}
    \item \texttt{file}: The Excel workbook.
    \item \texttt{department}: The target Department.
\end{itemize}

The importer automatically recognizes
two supported workbook formats.

\subsection{Platform Workbook Format}

The platform-specific workbook contains
three worksheets:

\begin{enumerate}
    \item Physical Devices.
    \item Network Elements.
    \item Measurements.
\end{enumerate}

References between records are represented
using \texttt{ref} and \texttt{@ref} columns.

\subsection{Client Workbook Format}

The importer also supports the workbook
format supplied by KIOS CoE, represented
by the file \texttt{EXCEL-317.xlsx}.

The workbook contains three principal
data worksheets.

\subsubsection{Measurements Worksheet}

The Measurements worksheet includes:

\begin{itemize}
    \item MeasurementId
    \item MeasurementName
    \item MeasurementCollectionBatch
    \item MeasurementType
    \item Meter Account
    \item Customer Account
    \item MeasurementUnit
    \item Details
    \item NetworkElementId
    \item PhysicalDeviceId
    \item DataCollectionType
    \item DataCollectionParameters
\end{itemize}

\subsubsection{NetworkElements Worksheet}

The NetworkElements worksheet includes:

\begin{itemize}
    \item NetworkElementId
    \item NetworkElementName
    \item NetworkElementType
    \item NetworkArea
    \item Latitude
    \item Longitude
    \item Elevation
    \item StrategicId
    \item DMAId
    \item Address
\end{itemize}

\subsubsection{Physical Devices Worksheet}

The \texttt{Physical\_devices} worksheet includes:

\begin{itemize}
    \item Device\_id
    \item Device\_name
    \item Device\_type
    \item CommunicationParameters
    \item longitude
    \item latitude
\end{itemize}

\subsection{Import Processing Rules}

The importer applies the following rules:

\begin{enumerate}

    \item Imported records are assigned to
    the Department selected in the import form.

    \item Any Department column contained in
    the workbook is ignored.

    \item The worksheets Statistics, LOVs,
    and Structure are ignored.

    \item Columns containing formulas are ignored.

    \item Relationships between imported records
    are resolved using the client's identifiers.

    \item Missing controlled-list values may be
    created only by System Administrators.

    \item If a required controlled-list value
    is missing and the user lacks permission
    to create it, the import is rejected.

    \item Controlled-list values longer than the
    permitted length are rejected with a
    column-level error.

    \item Each imported row is validated through
    the normal application serializers.

    \item Validation includes ownership rules
    and uniqueness constraints.

    \item The import operation is atomic.

\end{enumerate}

Atomic execution ensures that either all
valid records are committed or the entire
operation is rolled back.

\subsection{Validation and Error Reporting}

The importer reports validation errors
using the original workbook structure.

Error information identifies:

\begin{itemize}
    \item Worksheet name.
    \item Row number.
    \item Original client column name.
\end{itemize}

The following import restrictions apply:

\begin{table}[htbp]
\centering
\caption{Excel Import Restrictions}
\label{tab:phase3_import_limits}
\begin{tabularx}{\textwidth}{P{5cm} L}
\toprule
\textbf{Restriction} & \textbf{Limit} \\
\midrule
Maximum file size & 5 MiB \\
Maximum rows per worksheet & 2000 \\
Macro-enabled workbooks & Not supported \\
Import execution & Atomic transaction \\
\bottomrule
\end{tabularx}
\end{table}

\subsection{Data Collection Configuration Mapping}

The client workbook contains collection-related
fields that are stored in the
DataCollectionConfiguration entity.

The mapping includes:

\begin{itemize}
    \item \texttt{protocol}
    \item \texttt{extra\_settings}
    \item \texttt{source\_identifier}
    \item \texttt{collection\_batch}
\end{itemize}

This separates measurement metadata
from the configuration used for
data collection.

\subsection{Client Workbook Import Validation}

The actual client workbook was tested
against a test database.

The import completed successfully.

Table~\ref{tab:phase3_import_results}
summarizes the observed results.

\begin{table}[htbp]
\centering
\caption{Client Workbook Import Results}
\label{tab:phase3_import_results}
\begin{tabularx}{\textwidth}{L P{3cm}}
\toprule
\textbf{Imported Data} & \textbf{Result} \\
\midrule
Physical Devices & 11 \\
Network Elements & 10 \\
Measurements & 22 \\
New controlled-list entries & 28 \\
Approximate execution time & 0.5 seconds \\
\bottomrule
\end{tabularx}
\end{table}

These results demonstrate successful
processing of the supplied client workbook
within the test environment.

\subsection{Export Functionality}

The frontend supports exporting data
in CSV and XLSX formats.

Exported values are escaped to reduce
the risk of spreadsheet formula injection.

The exported files are generated
on the client side.

The export format is not intended
to serve as a template for re-importing
data into the platform.

% =========================================================
% 9. FRONTEND IMPLEMENTATION
% =========================================================

\section{Frontend Implementation}

\subsection{Frontend Architecture}

The frontend is implemented using React 19
and Vite.

It communicates with the Django backend
through the REST API.

The interface provides access to metadata
management, geographical visualization,
import/export operations, and administration
features.

\subsection{Implemented Pages}

Table~\ref{tab:phase3_frontend_pages}
summarizes the principal interface components.

\begin{longtable}{P{4cm} P{11cm}}
\caption{Frontend Pages and Functionality}
\label{tab:phase3_frontend_pages}\\

\toprule
\textbf{Page} & \textbf{Functionality} \\
\midrule
\endfirsthead

\toprule
\textbf{Page} & \textbf{Functionality} \\
\midrule
\endhead

\bottomrule
\endfoot

Dashboard &
Displays statistics, recent records,
and information grouped by area. \\

Map &
Displays physical devices and network
elements geographically. \\

Physical Devices &
Provides tables and forms for device management. \\

Network Elements &
Provides tables and forms for network element management. \\

Measurements &
Provides tables and forms for measurement management. \\

Controlled Lists &
Displays and manages reference data according
to the user's permissions. \\

Import/Export &
Supports Excel import and CSV/XLSX export. \\

API Keys &
Supports API key creation, one-time secret
display, and revocation. \\

Users &
Provides user management functionality. \\

Settings &
Displays role-specific configuration options. \\

Profile Menu &
Provides Department switching and logout. \\

\end{longtable}

\subsection{Dashboard}

The Dashboard provides an overview
of the metadata stored in the platform.

Its principal components include:

\begin{itemize}
    \item Statistical information.
    \item Recently created records.
    \item Information grouped by geographical area.
\end{itemize}

The information displayed depends on
the authenticated user's accessible scope.

\subsection{Interactive Map}

The Map page uses Leaflet together with a
MapLibre GL basemap provided by OpenFreeMap,
whose labels are displayed in Greek.

It displays geographical information
associated with physical devices
and network elements.

Geographical coordinates are represented
using the WGS84 coordinate system.

The map requires an internet connection
to retrieve the basemap style and tiles.

\subsection{Metadata Management Forms}

The interface provides tables and forms
for managing:

\begin{itemize}
    \item Physical Devices.
    \item Network Elements.
    \item Measurements.
    \item Controlled reference lists.
\end{itemize}

The available operations depend on
the user's authorization role.

\subsection{Role-Specific Interface}

The frontend adapts the Settings
interface according to the user's role.

For example, administrative operations
are displayed only when the corresponding
capabilities are available.

However, hiding interface elements
does not constitute authorization.

All permissions are independently
enforced by the backend.

\subsection{Department Switching}

The profile menu includes a
Department switcher.

When the active Department or
authorization role changes,
the workspace reloads its data.

This ensures that the interface
reflects the currently selected
organizational context.

% =========================================================
% 10. DEPLOYMENT AND EXECUTION
% =========================================================

\section{Installation and Execution}

\subsection{Environment Configuration}

The application uses environment variables
for deployment-specific configuration.

The example environment file is copied
to create the local configuration:

\begin{lstlisting}[language=bash]
cp .env.example .env
\end{lstlisting}

The environment configuration must contain
unique credentials and a
\texttt{DJANGO\_SECRET\_KEY}.

\subsection{Docker Compose Execution}

The application services are built
and started using Docker Compose.

\begin{lstlisting}[language=bash]
docker compose build
docker compose up -d
\end{lstlisting}

The deployment contains the database,
backend, and frontend services.

\subsection{Database Initialization}

Database migrations must be applied
before using the application.

The Platform Department must then
be provisioned using the Django
management command:

\begin{lstlisting}
provision_platform_department
\end{lstlisting}

The resulting Department identifier
is assigned to:

\begin{lstlisting}
METADATA_PLATFORM_DEPARTMENT_ID
\end{lstlisting}

\subsection{Access Policy Configuration}

Department-based authorization is
activated through the environment setting:

\begin{lstlisting}
METADATA_ACCESS_POLICY=platform_security.access.DepartmentAccessPolicy
\end{lstlisting}

Without the configured policy,
the default NoAccessPolicy denies
metadata operations.

\subsection{User Configuration}

Platform users are configured through
the Django management command:

\begin{lstlisting}
configure_platform_user
\end{lstlisting}

Passwords are provided interactively
rather than embedded in command arguments.

\subsection{Demo Data}

The project includes a management command
for generating demonstration data:

\begin{lstlisting}
seed_demo_data
\end{lstlisting}

This command supports testing
and demonstration of the application.

\subsection{Application Access}

After successful initialization,
the frontend is accessible at:

\begin{lstlisting}
http://localhost:8080
\end{lstlisting}

% =========================================================
% 11. TESTING AND VALIDATION
% =========================================================

\section{Testing and Validation}

\subsection{Testing Strategy}

The implementation was evaluated continuously during development
using automated backend, frontend, and end-to-end tests.
Testing focused on the correctness of metadata operations,
authorization rules, database migrations, import functionality,
and frontend behavior.

Acceptance, negative, and boundary testing of the final release
are outside the scope of this phase and are described in Phase~4.

\subsection{Backend Testing}

The backend test suite contains 591 pytest tests and 20 subtests.
The tests execute against a PostgreSQL/PostGIS database.

The covered areas include:

\begin{itemize}
    \item REST API behavior and error contracts.
    \item Authentication, sessions, login throttling, and expiry.
    \item Role-based permissions and Department-based access restrictions.
    \item API key creation, expiration, revocation, and scope.
    \item Database migrations, including upgrade and reverse paths.
    \item Excel import processing, including the client workbook format.
    \item Input validation and security-related behavior.
\end{itemize}

\subsection{Frontend Testing}

The frontend includes 89 unit and component tests in 18 test files,
implemented using Vitest.
In addition, ESLint is used for static analysis, and the production
build is verified with Vite.

\subsection{End-to-End Testing}

Three Playwright workflows run against the integrated application
(React frontend, Django backend, and PostGIS database).
They cover login and logout, Department switching, relational
create/read/update/delete operations, persistence, geographical
display, export, API key access, role restrictions, and
cross-Organisation isolation.

\subsection{Continuous Integration}

All checks are executed automatically by GitHub Actions on every push
and pull request. The workflow contains the jobs shown in
Table~\ref{tab:phase3_ci}.

\begin{table}[htbp]
\centering
\caption{Continuous Integration Jobs and Latest Results}
\label{tab:phase3_ci}
\small
\begin{tabularx}{\textwidth}{P{4.2cm} P{3.2cm} L}
\toprule
\textbf{Job} & \textbf{Tool} & \textbf{Result} \\
\midrule
Backend tests &
pytest, PostGIS &
591 passed, 20 subtests passed \\

Frontend tests &
Vitest &
89 passed (18 files) \\

Static analysis and build &
ESLint, Vite &
Passed \\

Integration validation &
Playwright, disposable PostGIS &
3 passed \\

Migration consistency &
Django &
No missing migrations \\
\bottomrule
\end{tabularx}
\end{table}

The integration job runs the full application in a fresh, disposable
PostGIS container through the script \texttt{validate\_disposable.sh},
so that no persistent database is accessed.

\subsection{Defects and Known Gaps}

Table~\ref{tab:phase3_bugs} lists the significant defects and gaps that
were identified during development, together with their status and
how they were addressed. Defects found by the automated tests or during
quality assurance were corrected before delivery.

{\small
\begin{longtable}{P{1.2cm} P{5.4cm} P{1.8cm} P{6.6cm}}
\caption{Significant Defects and Known Gaps}
\label{tab:phase3_bugs}\\

\toprule
\textbf{ID} & \textbf{Description} & \textbf{Status} & \textbf{Resolution} \\
\midrule
\endfirsthead

\toprule
\textbf{ID} & \textbf{Description} & \textbf{Status} & \textbf{Resolution} \\
\midrule
\endhead

\bottomrule
\endfoot

B01 &
External identifiers were unique across the whole platform instead of
within a Department. &
Fixed &
Uniqueness is scoped by Department, with a migration and regression
tests. \\

B02 &
The migration that scopes identifiers by Department failed when
upgrading an existing database. &
Fixed &
A reproducing test was added and the migration was rewritten to upgrade
safely. \\

B03 &
Deeply nested JSON input caused an unhandled decoder error. &
Fixed &
The error is handled as a parse error (HTTP~400) and the stored data
is preserved. \\

B04 &
A corrupt compressed stream inside an uploaded XLSX file was not
rejected cleanly. &
Fixed &
The file is rejected with a validation error. \\

B05 &
Nested external identifiers were not validated when a record was
updated with PATCH. &
Fixed &
Validation was added and covered by tests. \\

B06 &
A user who could read a record but not modify it was not clearly
distinguished from a user outside the record's scope. &
Fixed &
Forbidden writes return HTTP~403, while out-of-scope records return
HTTP~404. \\

B07 &
Exported CSV cells starting with formula characters could be evaluated
by spreadsheet applications. &
Fixed &
Such values are neutralized when the file is generated. \\

B08 &
Map selection could confuse a device and a network element with the
same numeric identifier. &
Fixed &
Selection keys are qualified by the resource type. \\

B09 &
The frontend could display stale or incomplete data when a request
failed or was superseded. &
Fixed &
Resources load independently, stale requests are cancelled, and
partial failures are reported. \\

B10 &
The Platform Department could be assumed implicitly. &
Fixed &
The Platform Department must be configured explicitly. \\

B11 &
A controlled-list name longer than the database limit caused an
HTTP~500 error during the client workbook import. &
Fixed &
The import reports a column-level validation error, with a regression
test. \\

B12 &
Two end-to-end tests failed after the Dashboard and the profile menu
were redesigned. &
Fixed &
The tests were updated to match the new interface. \\

B13 &
Historical migration tests assumed the previous database schema after
new migrations were added. &
Fixed &
The tests were updated to account for the new migrations. \\

G01 &
The frontend forms did not expose \texttt{network\_area},
\texttt{elevation}, \texttt{address}, and
\texttt{communication\_parameters}. &
Fixed &
The fields were added to the forms and the detail views, with input
validation and tests for creating, editing, and clearing them. \\

G02 &
Fields of the client's design workbook (sensor stations, communication
profiles, credentials, API endpoint, SQL query settings) are not
modelled. &
Open &
Documented as future work. \\

G03 &
The Docker Compose deployment is not hardened for public internet
access. &
Open &
Documented as a limitation; TLS termination and trusted client address
handling are required. \\

G04 &
Export is generated in the browser and limited to 100\,000 records. &
Open &
Documented as a known limitation. \\

G05 &
The map requires internet access to load its basemap. &
Open &
Documented as a known limitation. \\

\end{longtable}
}

\subsection{Testing Summary}

\begin{table}[htbp]
\centering
\caption{Implementation Testing Summary}
\label{tab:phase3_testing}
\small
\begin{tabularx}{\textwidth}{P{4cm} P{3cm} L}
\toprule
\textbf{Test Category} & \textbf{Tool} & \textbf{Result in CI} \\
\midrule
Backend tests & pytest & 591 passed \\
Frontend unit tests & Vitest & 89 passed \\
Static analysis & ESLint & Passed \\
Production build & Vite & Passed \\
End-to-end tests & Playwright & 3 passed \\
Isolated validation & Shell script & Passed \\
\bottomrule
\end{tabularx}
\end{table}

% =========================================================
% 12. IMPLEMENTATION LIMITATIONS
% =========================================================

\section{Implementation Limitations}

\subsection{Deployment Security}

The current Docker Compose configuration
is intended for local deployment.

It does not constitute a fully hardened
production deployment for public
internet access.

Additional configuration is required
for aspects such as:

\begin{itemize}
    \item TLS termination through a reverse proxy.
    \item Trusted client address handling.
\end{itemize}

\subsection{Unimplemented Client Fields}

Some fields present in the client's
design workbook are not yet represented
in the implemented domain model.

These include:

\begin{itemize}
    \item SensorStations.
    \item Communication\_Profiles.
    \item Credentials-related fields.
    \item API\_Endpoint.
    \item SQL\_Query-related fields.
\end{itemize}

These elements remain outside
the current implementation.

\subsection{Export Scalability}

CSV and XLSX export operations
are performed in the frontend.

The implementation retrieves
data into the client application
before generating the export file.

The export mechanism is limited
to 100,000 records and is not
designed for very large datasets.

\subsection{External Map Dependencies}

The interactive map uses a basemap
provided by OpenFreeMap.

Retrieving the map style and tiles requires
internet connectivity.

Offline tile availability is
not provided by the described
implementation.

% =========================================================
% 13. FUTURE WORK
% =========================================================

\section{Future Improvements}

The implementation provides a foundation
for further development.

Potential improvements include:

\begin{enumerate}

    \item Completing production deployment
    hardening, including TLS proxy
    configuration and trusted client
    address handling.

    \item Extending the domain model
    to support additional fields
    from the client's design workbook.

    \item Improving export scalability
    for larger datasets.

    \item Investigating approaches
    for map operation when external
    tile services are unavailable.

\end{enumerate}

% =========================================================
% 14. CONCLUSION
% =========================================================

\section{Conclusion}

Phase 3 focused on transforming
the requirements and architectural
design of the Sensor Metadata Management
Platform into an integrated software
implementation.

The resulting system combines
Django REST Framework, PostgreSQL/PostGIS,
React, and Docker Compose to support
centralized management of sensor metadata.

The implementation includes
Department-based authorization,
session and API key authentication,
REST API operations, geographical
visualization, and Excel import/export
functionality.

The successful import of the
client-provided workbook demonstrates
the ability of the system to process
real metadata using the implemented
validation and relationship rules.

Automated backend, frontend, and
end-to-end tests, executed through
continuous integration, provide
additional validation of the
implemented functionality.

Nevertheless, several limitations
remain, including production deployment
hardening and the client design fields
that are not yet modelled.

Overall, Phase 3 establishes an
integrated implementation of the
Sensor Metadata Management Platform
and provides a technical foundation
for subsequent improvements and
final system validation.

\end{document}
